\section{Conclusion}\label{sec:conclusion}

We have presented the first geometric analysis of loss landscapes in neural
theorem proving, comparing chain-of-thought decomposition with direct proof
generation. Our main theoretical contribution,
Theorem~\ref{thm:hessian-bound}, shows that the additive structure of the
decomposed loss bounds its Hessian spectral norm by the maximum per-step
spectral norm when per-step curvature directions are orthogonal. Our
experiments on propositional logic theorem proving with matched transformer
models confirm that CoT models converge to flatter minima with lower Hessian
top eigenvalue, lower SAM sharpness, and smaller gradient norms.

The most notable finding is the flat-but-isolated phenomenon: CoT minima are
individually wider but more separated from each other in parameter space, as
explained by the barrier superadditivity of
Proposition~\ref{prop:barrier-super}. This challenges the common assumption
that local flatness implies global connectivity and suggests that these are
independent geometric properties of the loss landscape.

Several directions merit further investigation.

\begin{enumerate}[leftmargin=*,itemsep=2pt]
\item \textbf{Scale.} Does the spectral orthogonality of per-step Hessians
  persist at the scale of modern theorem provers ($10^9$+ parameters)?

\item \textbf{Permutation alignment.} Applying Git
  Re-Basin~\cite{ainsworth2023git} before barrier measurement may reveal
  hidden connectivity between CoT solutions.

\item \textbf{SAM intervention.} Training direct models with
  sharpness-aware minimization~\cite{foret2021sharpness} would test whether
  flattening the landscape closes the performance gap.

\item \textbf{Decomposition depth.} Is there a monotone relationship between
  the number of CoT steps $K$ and the Hessian top eigenvalue? Our theory
  predicts diminishing returns once the per-step Hessians are approximately
  orthogonal, but this has not been empirically verified.

\item \textbf{Topological analysis.} Persistent homology of loss sublevel
  sets could characterize landscape topology beyond the barrier-based measures
  studied here.
\end{enumerate}

\begin{conjecture}\label{conj:phase}
  There exists a critical model capacity $d^*$ below which CoT and direct
  models converge to geometrically distinct regions of parameter space, and
  above which the landscape simplifies to a single basin (modulo permutation
  symmetries) for both strategies. The critical capacity $d^*$ depends on the
  complexity of the theorem proving domain.
\end{conjecture}

Understanding how proof decomposition shapes the optimization landscape is a
step toward principled design of neural reasoning systems. The geometric tools
developed here---spectral decomposition of the Hessian across reasoning steps,
barrier analysis for structured losses, and scale-dependent sharpness
comparison---are applicable beyond theorem proving to any setting where the
output can be decomposed into intermediate reasoning steps.
